% The four parts of the Day 3 lab. The steps follow Labs/day03/README.md.
%
% Hardware status: not tested on hardware. A value that only a board can
% give is written as "..." and the students measure it. The deck does not
% show a result that the students must predict.

% ------------------------------------------------------------------- Part A
\begin{frame}{Part A: train (45 min)}
  \small
  Start the notebook: \code{../../.venv/bin/jupyter lab
  motion\_classifier.ipynb}. It uses your recordings in
  \code{Labs/day02/data/}.
  \vskip 0.3em
  \renewcommand{\arraystretch}{1.25}
  \begin{tabular}{@{}L{0.09\textwidth}L{0.36\textwidth}L{0.45\textwidth}@{}}
    \toprule
    \textbf{Task} & \textbf{You write} & \textbf{The notebook checks} \\
    \midrule
    1 & The function that removes the mean and gives the RMS
      & A sine wave with the amplitude 1 gives an RMS of 0.707 \\
    2 & The two hidden layers of the model
      & The model has 1534 parameters \\
    3 & Your prediction: features against the raw window
      & Not checked. You compare it with the result. \\
    4 & The two lines of the converter (Part B)
      & The file starts correctly, and its output is equal to the output
        of Keras \\
    5 & Your estimate of the arena size (Part B)
      & Not checked. You measure it in Part C. \\
    \bottomrule
  \end{tabular}
  \vskip 0.5em
\end{frame}

\begin{frame}{Part A: features, model, and test}
  \begingroup\centering
  \begin{tikzpicture}[font=\footnotesize, >={Stealth[length=5pt]},
      b/.style={draw=coolgray, rounded corners=3pt, fill=boxgray,
                minimum width=1.75cm, minimum height=1.0cm, align=center},
      lab/.style={font=\scriptsize, text=coolgray, align=center}]
    \node[b] (a) at (0,0)    {Windows};
    \node[b, draw=cardinalred] (b) at (2.15,0) {Features};
    \node[b] (c) at (4.3,0)  {Normalize};
    \node[b] (d) at (6.45,0) {Model};
    \node[b] (e) at (8.6,0)  {Test};
    \foreach \a/\b in {a/b, b/c, c/d, d/e} { \draw[->] (\a) -- (\b); }
    \node[lab] at (0,-0.95)    {100 samples,\\3 axes};
    \node[lab] at (2.15,-0.95) {63 values};
    \node[lab] at (4.3,-0.95)  {mean and std of\\the training set};
    \node[lab] at (6.45,-0.95) {63, 20, 10, 4};
    \node[lab] at (8.6,-0.95)  {the test session\\of Day 2};
  \end{tikzpicture}\par\endgroup
  \vskip 0.5em
  \small
  \begin{itemize}
    \item \textbf{You see:} the plot of the spectral power for one window of
          each class. Find the axis and the frequency that separate your
          classes.
    \item \textbf{You record:} the training accuracy, the test accuracy, and
          the confusion matrix.
  \end{itemize}
  \begin{exerciseblock}
    \textbf{Prediction (task 3).} The notebook trains a second model with
    the raw window: 300 inputs. How many parameters does it have? Which
    input gives the higher test accuracy? Write the prediction before you
    run section 6.
  \end{exerciseblock}
\end{frame}

% ------------------------------------------------------------------- Part B
\begin{frame}[fragile]{Part B: convert, and write the files (50 min)}
  \small
  \textbf{Task 4.} Write the two lines of the converter in section 7 of the
  notebook. The notebook then prints lines of this form:
\begin{lstlisting}
Task 4: complete
LiteRT file:            ... bytes
operator types:         [...]
largest difference between Keras and LiteRT: ...
\end{lstlisting}
  \begin{itemize}
    \item The difference must be smaller than 0.00001.
    \item \textbf{You record:} the file size and the operator types. Task B1
          needs the operator types.
  \end{itemize}
  Section 8 writes three files into \code{sketches/motion\_classifier/}:
  \vskip 0.2em
  \renewcommand{\arraystretch}{1.2}
  \begin{tabular}{@{}L{0.26\textwidth}L{0.66\textwidth}@{}}
    \toprule
    \code{model.h} & The LiteRT file as a C array \\
    \code{model\_settings.h} & The class names, and the mean and the
        standard deviation of each feature \\
    \code{test\_window.h} & One test window, its 63 features, and the
        output of the model for it \\
    \bottomrule
  \end{tabular}
  \vskip 0.3em
  \textbf{Task 5.} Estimate the arena size from the table of the last cell.
  Write the estimate in the report before Part C.
\end{frame}

\begin{frame}{Part B: tasks B1 to B4, the inference code}
  \small
  Install the library Chirale\_TensorFlowLite 2.0.0. Open
  \code{sketches/motion\_classifier/motion\_classifier.ino}. Set
  \code{Tools > PSRAM > Disabled}.
  \vskip 0.3em
  \renewcommand{\arraystretch}{1.25}
  \begin{tabular}{@{}L{0.09\textwidth}L{0.24\textwidth}L{0.57\textwidth}@{}}
    \toprule
    \textbf{Task} & \textbf{Place} & \textbf{You write} \\
    \midrule
    B1 & \code{setupModel()} & The operator resolver: one \code{Add...()}
         call for each operator type of your model \\
    B2 & \code{setupModel()} & The interpreter, and the call of
         \code{AllocateTensors()} \\
    B3 & \code{classify()} & Normalize the features, write the input
         tensor, call \code{Invoke()}, read the output tensor \\
    B4 & \code{bestClass()} & The position of the largest probability \\
    \bottomrule
  \end{tabular}
  \vskip 0.5em
  \begin{itemize}
    \item The sketch compiles before you start. It then prints
          \code{Tasks B1 and B2 are not complete.}
    \item \textbf{You record:} the line \code{Sketch uses ... bytes} of this
          first build. It is the flash use with no runtime in the program.
  \end{itemize}
\end{frame}

\begin{frame}[fragile]{Part B: the self-test and the demonstration}
  \small
  Upload the complete sketch. The Serial Monitor at 115200 baud shows:
\begin{lstlisting}
Arena location:       internal RAM
Model size (flash):   ... bytes
Sketch size (flash):  ... bytes
Arena size:           4096 bytes
Arena used:           ... bytes
Self-test, features:  largest difference ...
Self-test, model:     largest difference ...
Self-test, class:     ... (expected: ...)
Self-test:            PASS
\end{lstlisting}
  \begin{itemize}
    \item The self-test calculates the features of the window of
          \code{test\_window.h} on the board, runs the model, and compares
          the results with the results of the notebook.
    \item \code{FAIL}: check task B3, and run section 8 of the notebook
          again.
    \item Then make each of the four motions. The display shows the class
          and its probability.
    \item \textbf{You record:} for each motion, the class that the board
          shows most of the time.
  \end{itemize}
  \source{The output shows the form of the lines. Nobody measured the
    values on the kit.}
\end{frame}

% ------------------------------------------------------------------- Part C
\begin{frame}[fragile]{Part C: measure three numbers (30 min)}
  \small
  \renewcommand{\arraystretch}{1.25}
  \begin{tabular}{@{}L{0.13\textwidth}L{0.40\textwidth}L{0.38\textwidth}@{}}
    \toprule
    \textbf{Number} & \textbf{Where you read it} & \textbf{You record} \\
    \midrule
    Flash   & \code{Sketch uses ... bytes} of the build, and
              \code{Model size} of the Serial Monitor
            & The two values, and the difference to the first build of
              Part B \\
    Arena   & \code{Arena used} of the Serial Monitor
            & The value, and your estimate of task 5 \\
    Latency & The columns \code{features\_us} and \code{invoke\_us}
            & The median and the maximum of 20 lines \\
    \bottomrule
  \end{tabular}
  \vskip 0.4em
  One line of the Serial Monitor for each inference:
\begin{lstlisting}
class,probability,features_us,invoke_us,late_samples
...,...,...,...,...
\end{lstlisting}
  \begin{itemize}
    \item The sum of the two times must be below 20\,000 microseconds: one
          sample period.
    \item \code{late\_samples} counts the samples that the loop read too
          late. Write the value after 20 lines.
  \end{itemize}
\end{frame}

\begin{frame}{Part C: the arena, and the PSRAM}
  \begin{columns}[T]
    \begin{column}{0.47\textwidth}
      \begin{block}{The arena size}
        \begin{enumerate}
          \item Set \code{kTensorArenaSize} to 1024. Upload. Write the
                message of the Serial Monitor in the report.
          \item Set it to the value \code{Arena used} plus 256 bytes.
                Upload. The self-test must pass.
          \item Compare the measured value with your estimate. Explain the
                difference in one sentence.
        \end{enumerate}
      \end{block}
    \end{column}
    \begin{column}{0.49\textwidth}
      \begin{block}{The second configuration}
        \begin{enumerate}
          \item Set \code{Tools > PSRAM > OPI PSRAM}.
          \item Change the line of the sketch to
                \code{\#define ARENA\_IN\_PSRAM 1}.
          \item Upload. The Serial Monitor prints \code{Arena location:
                PSRAM}.
          \item Write the same three numbers in the second column of the
                report.
        \end{enumerate}
      \end{block}
    \end{column}
  \end{columns}
  \vskip 0.4em
  \begin{exerciseblock}
    \textbf{Prediction.} Is the inference faster, slower, or equal when the
    arena is in the PSRAM? Write your answer and the reason before you
    measure.
  \end{exerciseblock}
\end{frame}

% ------------------------------------------------------------------- Part D
\begin{frame}{Part D: Edge Impulse Studio (25 min)}
  \small
  \begin{enumerate}
    \item \textbf{Impulse.} Open your project of Day 2. In \code{Create
          impulse}: window size 2000 ms, window increase 200 ms, the block
          \code{Spectral Analysis}, and the block \code{Classification}.
    \item \textbf{Features.} In \code{Spectral features}: FFT length 32.
          Generate the features.
    \item \textbf{Training.} In \code{Classifier}: two dense layers with 20
          and 10 neurons, 30 training cycles. Then run \code{Model testing}.
    \item \textbf{Library.} In \code{Deployment}: \code{Arduino library}.
          Add the ZIP file with \code{Sketch > Include Library > Add .ZIP
          Library}.
    \item \textbf{Sketch.} Open \code{sketches/ei\_motion\_inference/}.
          Change the first \code{\#include} line to the header of your
          library. Upload.
  \end{enumerate}
  \begin{itemize}
    \item \textbf{You record:} the test accuracy, and the estimates of the
          Studio for the inferencing time, the peak RAM, and the flash use.
    \item \textbf{You record:} the two times of the line \code{Predictions
          (DSP: ... ms, Classification: ... ms, ...)} and the line
          \code{Sketch uses ... bytes}.
    \item If the build of the library is not complete in time, use the
          estimates of the Studio.
  \end{itemize}
\end{frame}

\begin{frame}{Part D: the comparison table}
  \small
  \renewcommand{\arraystretch}{1.3}
  \begin{tabular}{@{}L{0.30\textwidth}L{0.30\textwidth}L{0.30\textwidth}@{}}
    \toprule
    \textbf{Number} & \textbf{Your sketch with TensorFlow Lite Micro}
    & \textbf{Library of Edge Impulse} \\
    \midrule
    Test accuracy & Section 5 of the notebook & \code{Model testing} \\
    Flash         & \code{Sketch uses ... bytes} & \code{Sketch uses ...
                    bytes}, or the estimate \\
    RAM           & \code{Arena used} & The peak RAM of the Studio \\
    Time for the features & \code{features\_us} & DSP time \\
    Time for the inference & \code{invoke\_us} & Classification time \\
    Kernels       & Reference kernels & Optimized kernels (ESP-NN) \\
    \bottomrule
  \end{tabular}
  \vskip 0.6em
  \begin{itemize}
    \item Write ``measured'' or ``estimate'' next to each number of the
          second column.
    \item The two paths use different units: microseconds and
          milliseconds. Write the unit.
  \end{itemize}
\end{frame}
